\section{方法}

\subsection{整体思路}
设 $\mathcal N$ 为一个带开放腿的张量网络，其开放标签集合记为 $\mathcal O$，内部收缩标签集合记为 $\mathcal I$。
将所有内部标签收缩后，可得到完整输出张量 $X \in \mathbb{R}^{d_1 \times \cdots \times d_{|\mathcal O|}}$.
相应地，端到端时间可写为
\[
T_{\mathrm{total}} = T_{\mathrm{contract}} + T_{\mathrm{emit}},
\]
其中 $T_{\mathrm{contract}}$ 表示生成输出之前的内部收缩、状态合并与近似压缩代价，$T_{\mathrm{emit}}$ 表示将结果写回完整输出张量的代价。
由于完整输出本身必须被生成，$T_{\mathrm{emit}}$ 往往难以避免，因此真正值得优化的主要是 $T_{\mathrm{contract}}$。
基于这一判断，本文把近似严格限制在内部传播的分隔状态上，而不直接压缩输出边界。

整条方法主线可以概括为四步：
\begin{enumerate}
    \item 将完整输出张量按若干开放输出轴划分为多个输出块，并按块执行收缩与写回；
    \item 在网络交互图上构造一棵划分树，把一般图上的收缩过程写成一系列合并；
    \item 对每个子树，用分隔状态来概括它对外仍然可见的信息；
    \item 沿树递归合并这些分隔状态，并在每次合并后重新压缩，以控制中间状态规模。
\end{enumerate}

\subsection{划分树与分隔状态}
为了统一组织一般图上的递归收缩，我们在张量网络的交互图上构造一棵划分树。
树上的每个节点 $v$ 对应一个子网络 $\mathcal N_v$；叶节点对应单个原始张量，内部节点对应两个子网络的合并，根节点对应整个网络。

对任意节点 $v$，我们关心三类标签：
\begin{itemize}
    \item \textbf{输出标签} $\mathcal O_v$：该子树最终仍通向输出侧的标签；
    \item \textbf{边界标签} $\mathcal B_v$：该子树与外部其余网络连接所需保留的标签；
    \item \textbf{切割标签} $\mathcal C_v$：当 $v$ 由两个孩子合并而成时，连接左右孩子并在本层被消去的标签。
\end{itemize}
根节点不再与外部网络相连，因此其边界标签为空；叶节点不再继续拆分，因此没有切割标签。

给定输出块 $\Omega_b$ 后，我们将节点 $v$ 的\emph{分隔状态}定义为
\[
M_v \in \mathbb{R}^{d_{\mathrm{open}}(v) \times d_{\mathrm{bnd}}(v)},
\]
其中
\[
d_{\mathrm{open}}(v)=\prod_{\ell\in\mathcal O_v} d_{\ell},
\qquad
d_{\mathrm{bnd}}(v)=\prod_{\ell\in\mathcal B_v} d_{\ell}.
\]
这里 $d_{\ell}$ 表示标签 $\ell$ 的维度，矩阵 $M_v$ 的行对应输出侧配置，列对应边界侧配置，因此它刻画了当前子树在输出块 $\Omega_b$ 下对父节点可见的局部行为。

直接显式存储 $M_v$ 往往代价过高，因此我们采用低秩因子形式
\[
M_v \approx A_v \bigl(B_v\bigr)^\top,
\]
其中
\[
A_v \in \mathbb{R}^{d_{\mathrm{open}}(v) \times r_v},
\qquad
B_v \in \mathbb{R}^{d_{\mathrm{bnd}}(v) \times r_v},
\]
$r_v$ 是节点 $v$ 的接口秩。
这种写法把输出侧与边界侧显式分开，也清楚地说明了本文的近似位置：输出边界保持不变，近似只发生在分隔状态上。

\subsection{树的递归合并}
在定义了分隔状态之后，整棵划分树上的计算就可以写成统一的递归。

\paragraph{叶节点。}
对叶节点 $v$，其子网络只包含一个原始张量，把张量重排成
\[
T_v \in \mathbb{R}^{d_{\mathrm{open}}(v) \times d_{\mathrm{bnd}}(v)}.
\]
若不做近似，则 $T_v$ 就是该叶节点的精确分隔状态；若做近似，则对 $T_v$ 做低秩压缩，得到
\[
M_v \approx A_v \bigl(B_v\bigr)^\top.
\]

\paragraph{内部节点。}
设内部节点 $v$ 的两个孩子为 $v_L$ 和 $v_R$，其分隔状态分别为
\[
M_{v_L} \approx A_{v_L} \bigl(B_{v_L}\bigr)^\top,
\qquad
M_{v_R} \approx A_{v_R} \bigl(B_{v_R}\bigr)^\top.
\]
父节点的构造可分为三步。

\paragraph{(i) 合并输出侧。}
父节点的输出标签满足
\[
\mathcal O_v = \mathcal O_{v_L} \cup \mathcal O_{v_R},
\qquad
\mathcal O_{v_L} \cap \mathcal O_{v_R} = \varnothing.
\]
因此，父节点的输出侧配置空间是左右孩子输出侧配置空间的笛卡尔积。
在因子层面，这一步对应于将两个孩子的输出侧因子做 Kronecker 积
\[
A_{\mathrm{mat}} = A_{v_L} \otimes A_{v_R},
\]
其列数为 $r_{v_L} r_{v_R}$。

\paragraph{(ii) 收缩切割标签，保留边界标签。}
设父节点保留的边界标签集合为 $\mathcal B_v$，本层被消去的切割标签集合为 $\mathcal C_v$。
则父节点未压缩的边界侧因子可以看成对两个孩子的边界侧因子沿 $\mathcal C_v$ 做收缩：
\[
B_{\mathrm{mat}} = B_{v_L} \times_{\mathcal C_v} B_{v_R}.
\]

\paragraph{(iii) 重新压缩。}
经过前两步后，父节点得到未压缩表示
\[
\widetilde M_v = A_{\mathrm{mat}} B_{\mathrm{mat}}^\top.
\]
由于其中的中间列数已经膨胀到 $r_{v_L} r_{v_R}$，若不立即重压缩，秩会在树上迅速增长。
因此，我们对 $\widetilde M_v$ 再做一次低秩压缩，写回统一的分隔状态形式
\[
M_v \approx A_v \bigl(B_v\bigr)^\top.
\]

\subsection{自适应秩选择}
在分隔状态递归中，每次叶节点构造和内部节点合并之后，都要对当前分隔状态重新压缩，以防止秩在树上不断膨胀。
最直接的做法是为所有节点统一指定一个固定秩 $r$。
这种方案实现简单，但它默认划分树上不同节点的分隔状态难度大致相同；这一假设在实际中通常并不成立。
不同节点对应的子树规模以及谱衰减特性都可能差别很大。
有些节点本身较简单，较低秩已经足够；另一些节点则包含更复杂的内部耦合，若仍采用统一秩预算，就会过早丢失关键信息。
因此，固定秩往往会同时带来两类问题：在简单节点上保留了多余通道，在困难节点上又压得不够。

基于这一观察，本文在每个待压缩节点上按局部误差容忍度自适应选择秩，而不是预先指定统一目标秩。
设某个待压缩接口矩阵为
\[
M \in \mathbb{R}^{m \times n},
\]
其奇异值分解为
\[
M = U \Sigma V^\top,
\qquad
\Sigma = \mathrm{diag}(\sigma_1,\sigma_2,\ldots).
\]
若取其秩-$r$ 截断近似
\[
M_r = U_r \Sigma_r V_r^\top,
\]
则 Frobenius 范数下的相对误差为
\[
\frac{\|M - M_r\|_F}{\|M\|_F}
=
\left(
\frac{\sum_{i>r} \sigma_i^2}{\sum_i \sigma_i^2}
\right)^{1/2}.
\]
因此，给定容忍度 $\tau$，我们选择满足
\[
\left(
\frac{\sum_{i>r} \sigma_i^2}{\sum_i \sigma_i^2}
\right)^{1/2}
\le \tau
\]
的最小秩 $r$。
在实现中，叶节点与内部合并节点可分别使用不同的容忍度，记为 $\tau_{\mathrm{leaf}}$ 与 $\tau_{\mathrm{merge}}$。

这种做法把“每个节点保留多少通道”转化成“每个节点允许损失多少相对能量”。
谱衰减快的简单节点会自动停在较低秩，而谱衰减慢的困难节点会自动保留更多通道。
因此，自适应秩选择的关键在于按局部难度重新分配秩预算。

\subsection{隐式草图合并}
仅靠自适应秩选择还不够。
在内部节点处，若左右孩子的秩分别为 $r_L$ 和 $r_R$，则父节点在重压缩之前自然会出现一个秩乘积 $r_L r_R$。
如果先显式构造合并后的大接口矩阵，再对它做普通低秩分解，就会遇到中间瓶颈。
为此，本文引入隐式草图合并：不显式形成完整的合并接口，而是通过算子访问直接恢复其主导子空间\cite{halko2011finding,heddes2026approximating}。

设合并后待压缩的父节点接口为
\[
\widetilde M \in \mathbb{R}^{m \times n}.
\]
在显式路径中，它可以写成
\[
\widetilde M = A_{\mathrm{mat}} B_{\mathrm{mat}}^\top,
\]
其中 $A_{\mathrm{mat}}$ 来自输出侧因子的组合，$B_{\mathrm{mat}}$ 来自边界侧因子沿切割标签的收缩。
隐式草图合并真正需要的不是 $\widetilde M$ 的全部条目，而是对线性算子 $\widetilde M x$ 和 $\widetilde M^\top y$ 的访问。
在本文的合并结构下，这两个乘法可写为
\[
\widetilde M x = A_{\mathrm{mat}} \bigl(B_{\mathrm{mat}}^\top x\bigr),
\qquad
\widetilde M^\top y = B_{\mathrm{mat}} \bigl(A_{\mathrm{mat}}^\top y\bigr).
\]
因此，整个压缩过程都可以建立在算子之上，而不必先显式形成完整的 $\widetilde M$。

在此基础上，我们借鉴 randomized SVD 的思想\cite{halko2011finding}。
首先生成高斯随机测试矩阵
\[
\Omega \in \mathbb{R}^{n \times \ell},
\qquad
\Omega_{ij} \sim \mathcal N(0,1),
\]
其中 $\ell = k + p$，$k$ 是目标秩的估计规模，$p$ 是过采样参数。
然后构造样本矩阵
\[
Y_0 = \widetilde M \Omega.
\]
这里的 $Y_0$ 并不是通过显式矩阵乘法得到，而是通过对 $\Omega$ 逐列或成块应用算子 $\widetilde M$ 得到。
当谱衰减较慢时，为了更稳定地提取主导子空间，还可以做 $q$ 次幂迭代：
\[
Y_{t+1} = \widetilde M \bigl(\widetilde M^\top Y_t\bigr),
\qquad t=0,1,\ldots,q-1.
\]
随后对最终样本矩阵做正交化，得到
\[
Q = \mathrm{orth}(Y_q) \in \mathbb{R}^{m \times \ell},
\]
其中 $Q$ 近似张成了 $\widetilde M$ 的主导列空间。
再将原问题压缩到小矩阵
\[
M_{\mathrm{proj}} = Q^\top \widetilde M 
\in \mathbb{R}^{\ell \times n}
\]
上，并对其做精确 SVD：
\[
M_{\mathrm{proj}} = \widetilde U \Sigma V^\top.
\]

到这一步为止，原来的大接口已经被转化成一个小矩阵上的谱截断问题。
直接采用固定秩可以直接保留前 $r$ 个奇异方向；而本文主线采用基于尾能量的自适应截断：给定合并容忍度 $\tau$，选择满足
\[
\left(
\frac{\sum_{i>r} \sigma_i^2}{\sum_i \sigma_i^2}
\right)^{1/2}
\le \tau
\]
的最小秩 $r$。
设据此得到的截断结果为
\[
\widetilde M \approx Q \widetilde U_r \Sigma_r V_r^\top,
\]
为了与全文统一的分隔状态表示保持一致，我们将其写回为因子形式
\[
A = (Q \widetilde U_r) \Sigma_r^{1/2},
\qquad
B = V_r \Sigma_r^{1/2},
\]
从而得到
\[
\widetilde M \approx A B^\top.
\]
如果说自适应秩选择决定“最终保留多少通道”，那么隐式草图合并决定“如何在不显式形成膨胀接口的情况下找到这些通道”。

\subsection{输出分块与块级状态复用}
完整输出张量通常不会一次性生成，而是沿若干开放输出轴划分为多个小块，按块逐步收缩并写回。
设全局开放标签顺序为
\[
\mathcal O = (\ell_1,\ell_2,\ldots,\ell_m).
\]
为了简单起见，本文选取其中前 $q$ 个开放标签作为分块轴，并将每个分块轴上的索引集合划分为若干局部片段。
一个输出块 $\Omega_b$ 即由这些局部片段的笛卡尔积给出。

这里的输出分块不只是写回时的调度单位。
更重要的是，在固定某个输出块后，部分开放标签在当前块中只剩下较小的局部自由度，因此各个子树所见到的输出侧规模也会随之改变。
这会直接影响分隔状态的行维大小，从而改变整个递归收缩过程中的中间状态规模。

进一步地，不同全局输出块在某些子树上往往诱导出相同的局部限制条件。
如果一个子树在两个不同输出块下看到的是同一组局部输出切片，那么该子树状态就可以跨块复用，而不必重复构造。
因此，本文在块级收缩过程中缓存已经构造过的分隔状态，并在后续输出块中按局部条件进行复用。
从系统角度看，输出分块与块间状态复用共同决定了完整输出计算的组织方式。

需要强调的是，为了保证比较公平，本文的精确基线与近似方法使用同一组输出块，并采用相同的块级写回接口；
它们的差别只在于内部收缩如何组织：精确基线对每个输出块直接执行精确全网络收缩，而本文方法则通过划分树上的分隔状态递归完成近似收缩。

\subsection{算法总览}
为了便于整体把握，\Cref{alg:astnc} 给出本文方法的主流程。
它把前文的几个核心部件统一起来：输出分块负责组织完整输出计算过程，划分树负责把一般图上的收缩写成递归合并，分隔状态负责在树上向上传播局部信息，而重压缩则用于控制中间状态规模。
划分树构造、局部块键缓存、精确收缩基线的实现接口以及随机化压缩的可重复性设置见附录~\Cref{app:method-details}。

\begin{algorithm}[t]
\caption{\ours{} 的完整输出计算主流程}
\label{alg:astnc}
\begin{algorithmic}[1]
\Require 张量网络 $\mathcal N$，输出块集合 $\{\Omega_b\}_{b=1}^B$，压缩配置 $\theta$
\Ensure 完整输出张量 $X$

\State 在 $\mathcal N$ 的交互图上构造划分树 $\mathcal T$
\State 初始化输出缓冲区 $X$
\State 初始化状态缓存 $\mathcal C$

\For{每个输出块 $\Omega_b$}
    \State $S_{\mathrm{root}} \gets \textsc{构造状态}(\mathrm{root}(\mathcal T), \Omega_b, \theta, \mathcal C)$
    \State $X[\Omega_b] \gets \textsc{根节点读出}(S_{\mathrm{root}})$
\EndFor
\State \Return $X$

\vspace{0.4em}
\Function{构造状态}{$v, \Omega_b, \theta, \mathcal C$}
    \If{$(v,\Omega_b,\theta)$ 对应状态已在缓存中可复用}
        \State \Return 缓存状态
    \EndIf

    \If{$v$ 是叶节点}
        \State 在输出块 $\Omega_b$ 的限制下形成局部张量 $T_v$
        \State 在 $(\mathcal O_v,\mathcal B_v)$ 上构造分隔状态 $M_v$
        \State 视需要对 $M_v$ 做压缩
        \State 将状态存入缓存并返回
    \Else
        \State $S_L \gets \textsc{构造状态}(v_L,\Omega_b,\theta,\mathcal C)$
        \State $S_R \gets \textsc{构造状态}(v_R,\Omega_b,\theta,\mathcal C)$
        \State 合并 $S_L$ 与 $S_R$ 的输出侧因子
        \State 沿切割标签 $\mathcal C_v$ 收缩边界侧因子
        \State 对合并后的分隔状态重新压缩
        \State 将状态存入缓存并返回
    \EndIf
\EndFunction

\vspace{0.4em}
\Function{根节点读出}{$S_{\mathrm{root}}$}
    \State 收回根状态的最后一条秩维
    \State \Return 与 $\Omega_b$ 对应的稠密输出块
\EndFunction
\end{algorithmic}
\end{algorithm}